% ======================================================================
\section{Why the remaining errors happen}\label{sec:analysis}

This section uses the logged decision trail of every clip, per-item error lists, and an audit of the gate (sources in
Appendix~\ref{app:sources}).

\subsection{How far could this pipeline go?}\label{sec:ceiling}
A ceiling analysis replaces one stage at a time by the correct answer from the labels, on the same detections, and
re-scores. It was run on a slightly older version of the system (25 hits and 27 wrong on the development set, 55 needed
sounds), because it needs saved pictures; the proportions, not the exact numbers, are the point.

\begin{table}[h]
\centering\small
\caption{One stage made perfect at a time (development set, slightly older version).}
\label{tab:ceiling}
\begin{tabular}{lrrr}
\toprule
stage made perfect & hits & wrong & cost (from 2.451) \\
\midrule
on-screen check & 30 & 18 & 1.915 \\
listeners (and their filters) & 34 & 24 & 1.859 \\
vetoes and filters & 30 & 27 & 2.169 \\
timing & 26 & 20 & 2.197 \\
\midrule
all stages perfect & 43 & 3 & 0.761 \\
\bottomrule
\end{tabular}
\end{table}

Two conclusions follow. First, the on-screen check and the listeners are the two largest losses. Second, even with every
stage perfect, 12 of 55 needed sounds are still missed: 5 are heard by no model at all (a hammer, an explosion, a clang,
two golf-club whacks) and 7 are too faint or too badly timed. This is the \emph{detector floor}: with today's detectors
no decision rule can find more than about four in five needed sounds. Every attempt to hear more of them found real
sounds -- 80\,\% of the extra sounds a listener proposed on the screening set were real -- but each new hit brought two to
four wrong pictures, because a newly heard sound must still pass the gate (Appendix~\ref{app:history}).

\subsection{The main limit: ``present'' is not ``making the sound''}\label{sec:presence}
The gate's most common error is to take an object being in the frame as that object making the sound now: a church
tower is visible, so the ringing bell is silenced although nothing shows it swinging; perched birds are taken as the
source of a call made off screen. In the gate audit this explains 23 of 36 gate errors. Two of the three votes
effectively test presence -- naming the object, and asking whether that object \emph{could} make the sound -- so the
majority tests presence too.

The second error is the opposite: the name vote answers ``nothing'' on 42\,\% of the stretches the labels call visible,
when the source is small, hidden, or its action is fine (a hand, footsteps of the person holding the camera). Those
sounds keep a picture that counts as an on-screen wrong picture.

Every attempt to fix the gate moved along one trade-off line. An object detector vote, a segmentation vote, a box-and-crop
check, an audio-visual synchrony model, a localisation model, a step-by-step ``name, find, zoom, ask'' chain and a
larger or different vision model each silenced more on-screen sounds only by silencing about as many needed ones -- two
wrong pictures removed (saving $2 \times 2 = 4$) for every needed sound lost (costing 4), so no gain. Moving from a
7-billion to a 27-billion-parameter vision model changed balanced accuracy only from 0.61 to 0.62. The question asked of
the frames, \emph{is this object making this sound right now?}, is about an action at an instant, and open vision models
verify presence well and action poorly. A much larger model (Qwen3-VL-235B) is the one untested option; it did not fit
the available hardware.

\subsection{How the question is asked matters}
Given options (a) and (b), the vision model picks (a) 88\,\% of the time in one test; asked yes/no it leans to ``no'';
asked for text, its answer can be cut off. The final system therefore reads how strongly the model prefers ``yes'' over
``no'' at the first word, minus the same for the opposite question. This removes the habits but not the weakness: as a
gate signal it reaches an AUROC of only 0.647 (0.5 is a coin flip, 1 is perfect).

\subsection{Are the labels too strict?}
In two blind checks, about 1 in 6 pictures counted as wrong was a real, wanted, off-screen sound the labels do not list,
and several others were real sounds with the source on screen. Absolute costs are therefore somewhat pessimistic; the
comparison between systems is not affected, because the same labels score all of them.

\subsection{A worked example}
Clip \fp{bell_miami} (development) has one needed sound: a bell from 0.2 to 14.5\,s, labelled not visible. The detectors
find it. The gate cuts the 14-s sound into three stretches. In each, the name vote answers ``church bell'' (it shares the
word \emph{bell}: visible), the a/b vote says not visibly happening, and the describe vote says visible. Two of three in
every stretch: the sound is silenced and the needed bell is missed. The frames show the church; they do not show the bell
ringing. The one vote that asked about the action was right and was outvoted by two votes that measure presence.
